\section*{Chapter \ref{ch:g12:animal-exchange} --- Animal Exchange Systems}

\begin{solution}{exo:g12:animal-exchange:1}
An amoeba is a single \omterm{def:g10:the-cell:cell}{cell}: every region is near the membrane. A large
animal is many \omterm{def:g10:the-cell:cell}{cell} layers thick; deep \omterm{def:g10:the-cell:cell}{cells} sit far beyond \omterm{def:g10:membranes-enzymes:diffusion}{diffusion}
range from the outer surface, so internal surfaces and \omterm{def:g12:animal-exchange:bulk}{bulk flow} are
required.
\end{solution}

\begin{solution}{exo:g12:animal-exchange:2}
Any three of: long length; circular folds; \omterm{prop:g12:animal-exchange:villi}{villi}; microvilli; thin
epithelium; rich capillary/lymphatic bed; moist surface.
\end{solution}

\begin{solution}{exo:g12:animal-exchange:3}
Open: \omterm{def:g12:animal-exchange:open}{haemolymph} leaves vessels into sinuses (many arthropods, most
molluscs). Closed: blood stays in vessels and capillaries (vertebrates,
annelids, cephalopods).
\end{solution}

\begin{solution}{exo:g12:animal-exchange:4}
\omterm{def:g12:animal-exchange:gas}{Gas exchange} is \omterm{def:g10:membranes-enzymes:diffusion}{diffusion} of respiratory gases across a surface; mainly
O$_2$ in and CO$_2$ out for aerobic animals.
\end{solution}

\begin{solution}{exo:g12:animal-exchange:5}
Air $\to$ airways $\to$ \omterm{def:g12:animal-exchange:organs}{alveoli} $\to$ blood (bind haemoglobin) $\to$
pulmonary veins $\to$ left heart $\to$ aorta $\to$ systemic arteries
$\to$ muscle capillary $\to$ interstitial fluid $\to$ \omterm{def:g10:the-cell:cell}{cell} $\to$
\omterm{def:g10:cellular-respiration:mito}{mitochondrion}.
\end{solution}

\begin{solution}{exo:g12:animal-exchange:6}
Opposite flow keeps blood meeting water that is always slightly richer
in O$_2$, so the gradient persists along the lamella; concurrent flow
equilibrates midway and extracts less.
\end{solution}

\begin{solution}{exo:g12:animal-exchange:7}
Tracheae deliver air near mitochondria. \omterm{def:g10:membranes-enzymes:diffusion}{Diffusion} path length grows with
size; dense branching and ventilation help, but pure \omterm{def:g10:membranes-enzymes:diffusion}{diffusion} limits
very large bodies especially if atmospheric O$_2$ is low.
\end{solution}

\begin{solution}{exo:g12:animal-exchange:8}
Fish: pressure drops across \omterm{def:g12:animal-exchange:organs}{gills}, so systemic flow is at lower
pressure. Mammals repressurise oxygenated blood in the left heart for
strong systemic flow without overloading \omterm{def:g12:animal-exchange:organs}{lungs} --- needed for high
metabolic demand on land.
\end{solution}

\begin{solution}{exo:g12:animal-exchange:9}
$A/V = 6L^2/L^3 = 6/L$. Doubling $L$ halves $A/V$.
\end{solution}

\begin{solution}{exo:g12:animal-exchange:10}
Large \omterm{def:g12:animal-exchange:organs}{gill} area (many lamellae), very thin epithelium, continuous
ventilation (ram or pumped), countercurrent blood flow, high cardiac
output and haemoglobin with appropriate affinity for cold hypoxic water.
\end{solution}

\begin{solution}{exo:g12:animal-exchange:11}
Many \omterm{def:g10:molecules-of-life:lipids}{lipids} travel via lacteals as chylomicrons into lymph, then blood;
glucose enters blood capillaries directly (portal to liver). Both need
\omterm{def:g12:animal-exchange:bulk}{bulk flow} to distribute products beyond the gut wall.
\end{solution}

\begin{solution}{exo:g12:animal-exchange:12}
Fluid thickens the \omterm{def:g10:membranes-enzymes:diffusion}{diffusion} barrier and reduces effective air-facing
area / gradients; flux falls even if breathing rate is unchanged.
\end{solution}

\begin{solution}{pb:g12:animal-exchange:1}
\textbf{Part I.} (1)~Flatworm thin: short paths; mammal thick: long
paths need \omterm{def:g12:animal-exchange:bulk}{bulk flow} and internal surfaces. (2)~$A/V$ falls by factor
$100$ when $L$ rises from $1$ to $100$~mm. (3)~\omterm{def:g12:animal-exchange:organs}{Alveoli}, \omterm{prop:g12:animal-exchange:villi}{villi}/microvilli,
capillary beds (or nephrons).

\textbf{Part II.} (4)~Regional specialisation and continuous one-way
processing. (5)~Epithelial uptake into capillaries/lymph; bulk removal
maintains gradients. (6)~Herbivores often longer/more chambered guts and
greater microbial surface time for \omterm{def:g10:the-cell:support}{cell walls}.

\textbf{Part III.} (7)~Closed. (8)~Pulmonary then systemic repressurised
circuits. (9)~\omterm{def:g12:animal-exchange:organs}{Gills}/water vs \omterm{def:g12:animal-exchange:organs}{lungs}/air; ventilation differs.
(10)~Countercurrent vs tidal alveolar design with huge area.

\textbf{Part IV.} (11)~All increase together; local arterioles dilate in
muscle. (12)~Matched design: e.g.\ alveolar area and cardiac output
scale with demand; compare tuna \omterm{def:g12:animal-exchange:organs}{gills} vs mammal \omterm{def:g12:animal-exchange:organs}{lungs}.
\end{solution}
